\documentclass[11pt,letterpaper]{article}
\usepackage[T1]{fontenc}
\usepackage{lmodern,microtype,amsmath,amssymb,amsthm,booktabs,array,geometry,xcolor,fancyhdr,titlesec}
\usepackage[hyphens]{url}
\usepackage[colorlinks=true,linkcolor=teal!65!black,urlcolor=teal!65!black,citecolor=teal!65!black]{hyperref}
\geometry{margin=.92in,headheight=15pt}
\definecolor{navy}{HTML}{183547}\definecolor{ac}{HTML}{116962}
\hypersetup{pdftitle={Geometric progressions and nonabelian zero tests after Dirichlet GRH}}
\titleformat{\section}{\Large\bfseries\sffamily\color{navy}}{\thesection}{.65em}{}
\titleformat{\subsection}{\large\bfseries\sffamily\color{navy}}{\thesubsection}{.65em}{}
\pagestyle{fancy}\fancyhf{}\lhead{\small\sffamily\color{navy}Methods from the Dirichlet GRH proof}\rhead{\small\sffamily Research note}\cfoot{\small\sffamily\thepage}
\setlength{\emergencystretch}{2em}\raggedbottom
\setlength{\parskip}{.25em}
\newtheorem{theorem}{Theorem}[section]
\newcommand{\src}[2]{\href{https://github.com/lsernau/grh-lean/blob/72fead102f663d5c7a407964a13dc0cf9d977207/grh/#1}{#2}}
\newcommand{\prior}[2]{\href{https://github.com/VladimirReshetnikov/ProveIt/blob/48efe71a726dddc5a505534edaa4c0fd190416b9/docs/incoming/#1}{#2}}
\begin{document}
\thispagestyle{plain}
{\sffamily\small\bfseries\color{ac}NUMBER THEORY AND COMBINATORICS\hfill 25 NOVEMBER 2026}\par\vspace{8pt}
\noindent{\sffamily\bfseries\fontsize{25}{28}\selectfont\color{navy} Geometric progressions and\\ nonabelian zero tests\par}
\noindent{\large\sffamily Three transfers from the Lean proof of Dirichlet GRH\par}
\vspace{4pt}\noindent{\color{ac}\rule{\linewidth}{1pt}}

\noindent The Lean development for Dirichlet $L$-functions contains a progression-free construction on exponents and a one-sided criterion detecting zeros. Following these ideas through the earlier ProveIt reports leads to answers to two published density questions and to finite-degree results for nonabelian zero tests.
\begin{itemize}
\item McNew asked whether sets of positive integers avoiding three-term geometric progressions have different extremal natural and upper densities. We prove $\alpha<\overline\alpha$ for rational ratios and $\beta<\overline\beta$ for integer ratios (Questions 7.1 and 7.3). The respective gaps exceed $0.002836$ and $0.069218$. The obstruction already occurs in $\{1,\ldots,12\}$.
\item For every finite field, rational-ratio progression-free subsets of $\mathbb F_q[x]$ have strictly different optimal natural and upper densities. Both extrema are attainable, and the natural-density extremum is a computable real. Over $\mathbb F_2$ the gap exceeds $0.003660$. A change affecting just three polynomials explains the strict inequality.
\item For the nonabelian Artin factor of the cubic $X^3-X-1$, the eta-product representation yields strict quadratic Tur\'an inequalities at every shift. Our finite-height certificate gives all Jensen degrees through $16$ at every shift. Using Ng's published zero count extends the range to $945{,}173$. An exact negative cubic discriminant for a different kernel disproves the general sufficient condition in Dimitrov--Lucas, Theorem~1. A moment argument identifies the allowable first-order perturbations of a Gaussian if every cubic is to remain real-rooted.
\item For $G=\mathrm{PSL}_2(\mathbb F_q)$ in odd characteristic, $q\ge5$, we determine the smallest degree of a pointwise-nonnegative genuine character containing all nontrivial irreducibles. We also compute the minimum degree of a single coset action containing every irreducible. The classification includes the four exceptional small fields and the elementary case $q=3$.
\end{itemize}
These are mathematical proofs with auxiliary integer computations. The analytic result for the cubic field leaves its full critical-line conjecture open. All nine accompanying checkers use Python integer or rational arithmetic.

\section{What the Lean proof makes available}
\noindent\src{src/Statements.lean}{The public definitions} concern the Riemann zeta function and complex Dirichlet characters of positive modulus. In particular \src{src/Main.lean}{the exported nonvanishing theorem} covers $\Re s>1/2$ for Dirichlet $L(s,\chi)$ with the principal pole excepted. The primitive and general critical-strip formulations are separate exports. The pinned repository reports builds and kernel checks for its readable and compact editions on Lean 4.35.0, with the axioms \texttt{propext}, \texttt{Classical.choice} and \texttt{Quot.sound}. This study used source inspection and the published checks; it did not perform a fresh build. The source revision is \texttt{72fead102f66}, with Mathlib at \texttt{56a005855ea1}. The full commit and file identities are recorded in \texttt{source\_provenance.json}.

\paragraph{Exponents.}
In \src{src/New/Reduction/Data/Existence.lean\#L80-L116}{\texttt{exists\_three\_term\_free}}, an integer encodes a short digit vector. The construction bounds all digits so a weighted congruence implies exact coordinate equalities. Restricting to a sphere then rules out a strict weighted midpoint. In the polynomial ring the prime valuations already give integral coordinates. A nonconstant rational-function ratio yields a nonzero signed step in those coordinates. This leads to a concrete issue absent from the Lean lemma: the degree cutoffs weight configurations near their boundary disproportionately. We prove that choices maximizing different cutoffs are incompatible.

\paragraph{Zeros.}
The \src{src/New/Reduction/MainInequality.lean}{covolume inequality} compares two estimates for the same lattice, using short spanning forms and determinants from twisted theta lifts. It gives a \src{src/New/Reduction/OneSidedBound.lean}{one-sided estimate} for a smoothed prime sum involving a congruence class and a quadratic twist. The \src{src/New/Reduction/Zeros/Criterion.lean}{zero criterion} chooses a twist to prevent cancellation of zero orders in a prescribed region and applies a Landau argument. Extending this strategy to a nonabelian field requires an arithmetic estimate for the new Frobenius weights. We examine two independent pieces that can already be decided: what a specific completed $L$-function's moments imply, and the least degree of a character that gives every irreducible a positive multiplicity.

\paragraph{Earlier results used for comparison.}
Asada, Fourakis, Manski, McNew, Miller and Moreland studied polynomial ratios and explicitly asked about rational ratios in Question~6.4. Their formulation uses upper density; the distinction between the natural and upper optima proved here is a further question. McNew developed finite-prime methods over the integers. His later divisor-graph work gives an asymptotic with error term for the maximum size in each initial interval; simultaneous optimization across initial intervals is the issue here. See \href{https://arxiv.org/abs/1808.04923}{his later Theorem 9}. See the \href{https://arxiv.org/pdf/1512.01932}{finite-field paper} and the \href{https://arxiv.org/abs/1310.2277}{integer paper}. The \prior{grh-progression-densities-and-cubic-prime-race.zip}{incoming progression and cubic report} gives an upper-density improvement over the greedy set and a conditional prime-race sign for the discriminant-$-23$ field. The \prior{grh-methods-polynomial-density.zip}{subsequent polynomial-density report} gives finite-fiber approximations to the upper optimum. Our natural-density argument needs one compatible rule at every cutoff. The \prior{psl2_grh_transfer.zip}{character report} studies even characteristic and asks for the odd-characteristic degree optimization. Sections below give the resulting distinctions without assuming GRH for nonabelian Artin functions. Historical priority for the present conclusions has not been established.
\section{Two density questions of McNew over the integers}
Nathan McNew asked whether $\alpha<\overline\alpha$ and whether $\beta<\overline\beta$ in Questions 7.1 and 7.3 of \emph{On sets of integers which contain no three terms in geometric progression}, Math. Comp. 84 (2015), 2893--2910 (\href{https://www.nathanmcnew.com/GPFsets.pdf}{author's manuscript, section 7}). Here $\alpha$ is the supremum of natural densities of subsets of the positive integers avoiding three-term geometric progressions with rational ratio $r>1$. The barred quantity permits upper density. The pair $\beta,\overline\beta$ uses integer ratios $r>1$. The comparison below answers both questions. The proof uses unique factorization and a comparison of two optimizations over the same factors. It is independent of RH and GRH.

\begin{theorem}\label{th:integer-gap}
For both ratio conventions, the supremum of upper densities exceeds the supremum of natural densities by more than $0.000404$.
\end{theorem}

\paragraph{Two optimizations.}
Fix one of the ratio conventions and an integer $Y\ge 12$. Call an admissible set one containing no forbidden progression. Write $S$ for the positive integers composed solely of primes $p\le Y$, including $1$. Every positive integer has a unique expression $gh$ with $h\in S$ and $g$ divisible only by primes $p>Y$. Put
\[
 c=\prod_{p\le Y}(1-p^{-1}),\quad v=\prod_{p>Y}(1-p^{-2}),\quad
 B=c\sup_{A\subset S\ {\rm admissible}}\sum_{h\in A}\frac1h,\qquad
 U=c\int_1^\infty J(t)\frac{dt}{t^2},
\]
where $J(t)$ is the maximum size of an admissible subset of $S\cap[1,t]$. Writing $H(t)=\#(S\cap[1,t])$, Tonelli gives $c\int_1^\infty H(t)t^{-2}dt=c\sum_{h\in S}1/h=1$, hence $U\le1$. McNew's Theorem 4.2 uses finite smooth factors to study upper density. We need the additional comparison with one fixed admissible choice across all cutoffs.

\paragraph{Upper bound for natural density.}
Suppose $E$ is admissible and has natural density $d$. Partial summation gives $\sum_{n\le X,\ n\in E}1/n=(d+o(1))\log X$. For each rough $g$ the fixed fiber $A_g=\{h\in S:gh\in E\}$ is admissible. Consequently
\[
 \sum_{n\le X,\ n\in E}\frac1n\le\frac{B}{c}
           \sum_{g\le X,\ g\ {\rm rough}}\frac1g=(B+o(1))\log X.
\]
The last estimate follows by periodicity of the rough condition modulo $\prod_{p\le Y}p$ and partial summation. Thus $d\le B$.

\paragraph{Lower bound for upper density.}
The squarefree rough integers have density $cv$. To see this, first exclude the small primes and the squares of primes through a fixed bound by the Chinese remainder theorem, then use the upper density bound $\sum_{p>K}p^{-2}$ for the omitted squares, obtained by summing $\lfloor X/p^2\rfloor$. In any progression composed only of integers with squarefree rough parts, all valuations at primes $p>Y$ form arithmetic progressions in $\{0,1\}$. Their equality forces a common rough part. Thus it is sufficient to impose admissibility independently on each such fiber.

Choose integers $X_j$ with $X_j/X_{j-1}\to\infty$, starting at $X_0=0$. For each squarefree rough $g$ in $(X_{j-1},X_j]$ permanently use a maximizing subset of $S\cap[1,X_j/g]$. Their union over all $g,j$, after multiplying each subset by its $g$, is admissible. At $X=X_j$, fix $K>1$. For all sufficiently large $j$, the newly assigned rough parts in $(X/K,X]$ alone give, after division by $X$, a contribution tending to
\[
 cv\int_1^K J(t)\frac{dt}{t^2}.
\]
Indeed, $J$ is a step function with finitely many jumps on $[1,K]$. On each constant interval $(a,b)$ the corresponding rough parts in $(X/b,X/a)$ have normalized count tending to $cv(1/a-1/b)$; endpoints contribute $o(1)$. Let $K\to\infty$. The upper density of the constructed set is at least $vU$.

\paragraph{The incompatible choices through 12.}
An admissible subset maximizing its cardinality at every integer cutoff $k\le12$ would have to agree with the increasing greedy construction through $12$. Inductively, whenever greedy takes $k$, adding it to the agreed prefix improves the count, hence every simultaneous maximizer must take it. Whenever greedy rejects $k$, taking it would violate admissibility. All integers through $12$ belong to $S$. Under both conventions the greedy prefix is
\[
 \{1,2,3,5,6,7,8,10,11\}.
\]
The following admissible set has ten elements:
\[
 \{1,2,5,6,7,8,9,10,11,12\}.
\]
Indeed the complete list of rational-ratio progressions through $12$ is
$(1,2,4)$, $(1,3,9)$, $(2,4,8)$, $(3,6,12)$ and $(4,6,9)$. The first four have integer ratios. To enumerate them, use the form $(ab^2,abc,ac^2)$ with $1\le b<c$, $\gcd(b,c)=1$, and $ac^2\le12$.

For each admissible $A\subset S$ its prefix count therefore falls short of $J(k)$ by at least one at some $k\le12$. The loss persists throughout $[k,k+1)$ and contributes at least $1/(12\cdot13)$ to the corresponding integral. Since $\int_1^\infty |A\cap[1,t]|t^{-2}dt=\sum_{h\in A}1/h$, this proves $U-B\ge c/156$. With $u$ and $b$ denoting the extremal upper and natural densities,
\[
 u-b\ge vU-B\ge \frac{c}{156}-(1-v)
       \ge \frac{1}{156}\prod_{p\le Y}(1-p^{-1}) - \sum_{p>Y}p^{-2}.
\]
In particular Mertens' product theorem and $\sum_{n>Y}n^{-2}\le1/Y$ already give strict positivity for sufficiently large $Y$.

\paragraph{Explicit bound.}
Taking $Y=2000$, enumerate primes through $100000$ by Eratosthenes' sieve and bound the residual sum by $1/100000$. Directed fixed-point integer arithmetic in \texttt{check\_integer\_gap.py} verifies
\[
 \frac{1}{156}\prod_{p\le2000}(1-p^{-1})-
 \sum_{2000<p\le100000}p^{-2}-\frac{1}{100000}
 >\frac{404}{10^6}.
\]
The checker independently enumerates the prefix optimizers under both conventions. This proves Theorem~\ref{th:integer-gap}. The argument leaves McNew's Conjecture 7.2, on the exact value of $\alpha$, undecided.

\paragraph{A stronger numerical separation for integer ratios.}
For this convention the absolute lower bound $\overline\beta>0.818410$ of McNew's Theorem 4.4 allows a separate calculation. Use smooth primes $2,3,5,7$ in the upper bound for natural density, with $c=8/35$. There are $265$ such smooth numbers through $5000$. The maximum reciprocal sum among their subsets avoiding integer-ratio progressions is
\[
 M=\frac{1619153731921}{497871360000}.
\]
Assigning the full remaining smooth mass to the untested tail gives
\[
 \beta\le 1-c\left(\sum_{\substack{h\le5000\\ h\ \mathrm{smooth}}}\frac1h-M\right)
 =\frac{5245322990569}{7001316000000}<0.749192.
\]
It follows that $\overline\beta-\beta>0.069218$. The exhaustive calculation uses the common denominator of all $265$ reciprocal weights. In ascending order branch on including or excluding the least undecided integer. Whenever the first two terms of a forbidden progression have been selected, remove its third term from the remaining candidates. The selected and remaining weights together bound the optimum on that branch. The checker uses integer arithmetic for all pruning decisions.

\subsection{A larger separation from the first eighty integers}
The twelve-integer proof gives a common lower bound using a single incompatible prefix. We can count the loss over all prefixes through $N$. For $Y\ge N$, all those integers are $Y$-smooth. With $J(k)$ as above, put
\[
 L_N=\sum_{k=1}^N\frac{J(k)}{k(k+1)}
      -\max_{A\subseteq[1,N]\ {\rm admissible}}
              \sum_{a\in A}\left(\frac1a-\frac1{N+1}\right).
\]
The second expression integrates $|A\cap[1,t]|/t^2$ from $1$ to $N+1$.
The remaining integral of any prefix deficit is nonnegative. Thus $U-B\ge cL_N$. For $N=80$ the exact values are
\[
 \begin{array}{lrrl}
 \text{ratio}&L_{80}&Y&\text{certified lower bound on }u-b\\\hline
 \text{rational}&1553/45360&293&0.002836\\
 \text{integer}&16567/132300&83&0.013336 .
 \end{array}
\]
To reproduce the finite maximum, form each forbidden triple on $[1,k]$ as an edge with three vertices. Maximizing retained weight equals the total weight minus the minimum weight of a set meeting every edge. If an uncovered edge is $\{i,j,l\}$, branch on including each of its three vertices in the cover and remove all incident edges. The optimum is the minimum over these branches, with zero cost for the empty edge set. The program caches identical remaining sets of edges. It applies the same recursion with unit weights for each $J(k)$ and with integer multiples of the displayed rational weights for the second term. It then lower bounds $cL_{80}-\sum_{p>Y}p^{-2}$ with the same prime cutoff $100000$. This is \texttt{check\_integer\_prefix.py}. No floating-point optimizer is used.

\subsection{Computability and attainment of the natural-density optima}
For each ratio convention the natural-density supremum $b$ is a computable real and there is an admissible set with natural density $b$.
First, $vB\le b\le B$. For the lower bound choose a finite set $A$ approaching the weighted optimum defining $B$, and use this same fiber at every squarefree rough $g$. Its density is $cv\sum_{a\in A}1/a$, by the density of the permitted rough integers and a finite sum. Since $B\le1$, replacing $b$ by $B$ costs at most $1/Y$. For a fixed $Y$ the finite weighted optimization through $H$ costs at most $1-c\sum_{h\in S,\ h\le H}1/h$, which tends to zero and is an explicitly computable rational. This gives both-sided effective bounds without identifying the exact optimizer.

Here are details of attainment. Write $D(n)$ for the largest prime dividing $n$, with $D(1)=1$. Let $F_x(y)$ be the proportion of positive integers at most $x$ with $D(n)\le y$. For every fixed positive integer $a$,
\[
                 \sup_{y\ge1}|F_{x/a}(y)-F_x(y)|\longrightarrow0 .
\]
To see this, Dickman's theorem gives convergence of the distribution functions of $\log D(n)/\log x$ to the continuous law with distribution function $t\mapsto\rho(1/t)$ on $(0,1]$, extended by zero on $t\le0$ and one on $t\ge1$. Pointwise convergence of distribution functions to a continuous distribution function is uniform, by choosing a finite grid of quantiles and using monotonicity. Replacing $\log x$ by $\log(x/a)$ rescales $t$ by a factor tending to one, with uniform error tending to zero by continuity on a compact interval and the constant tails. This proves the assertion. See Dickman, \emph{On the frequency of numbers containing prime factors of a certain relative magnitude} (1930), and \href{https://library.slmath.org/books/Book44/files/09andrew.pdf}{Granville, \emph{Smooth numbers}, section 1}.

Choose cutoffs $Y_j\to\infty$ and finite smooth sets as above so the product rules $E_j$ have densities $e_j\to b$. For $\epsilon_j=2^{-j}$ each rule is approximable, in proportion on $[1,x]$ uniformly for integer $x\ge1$, within $\epsilon_j$ by a finite Boolean combination $C_j$ of divisibility conditions. Indeed membership in a finite list of exact smooth parts is a finite divisibility condition, and ignoring prime squares beyond $K$ costs at most $\sum_{p>K}p^{-2}$. Inclusion--exclusion writes the indicator of $C_j$ as a finite linear combination of $\mathbf1_{a\mid n}$. For $y\ge D(a)$, the joint probability of $a\mid n$ and $D(n)\le y$ is $\lfloor x/a\rfloor F_{x/a}(y)/\lfloor x\rfloor$. The uniform estimate just proved shows that this equals $a^{-1}F_x(y)+o_a(1)$ as $x\to\infty$, uniformly in such $y$.

Choose increasing thresholds $m_j>Y_j$ large enough that, for every $x,y\ge m_j$, the joint probability for $C_j$ and $D\le y$ differs from $d(C_j)F_x(y)$ by at most $\epsilon_j$. Using $E_j$ everywhere instead costs at most $2\epsilon_j$ more. On the bin $m_j<D(n)\le m_{j+1}$ use rule $E_j$. Omit the initial bin. In any forbidden triple, its largest prime would occur to exponent at most one wherever present, by squarefreeness above the local cutoff. The arithmetic progression of its three valuations must then be constant. All three members follow the same rule, a contradiction. For each bin the probability retained differs from $e_j$ times its probability by at most $6\epsilon_j$ (and both are zero if $x<m_j$). The mass of finitely many initial bins tends to zero since $F_x(m)\to0$ for fixed $m$. The errors on indices $j\ge J$ sum to at most $6\sum_{j\ge J}\epsilon_j$. Letting $x$ and then $J$ tend to infinity proves that this single set has natural density $b$.

The same weighted upper bound applies to every admissible set if one takes the limsup of its logarithmic means. Since an extremal set of natural density $b$ realizes logarithmic density $b$ as well, the supremum of upper logarithmic densities equals $b$. Rankin's stronger optimality conjecture therefore also has a finite-witness formulation. A finite, admissible set on some $Y$-smooth integers with reciprocal sum exceeding the full $Y$-smooth greedy sum would refute it: keep the greedy exponent rule at all outside primes and use the new local choice. The three outside valuation vectors of a progression would have to be equal. The converse follows because the normalized local optima converge to $b$ and the normalized local greedy sums converge to Rankin's density.

\section{The largest natural density for rational geometric progressions}
Fix a prime power $q$. We work with monic polynomials over $\mathbb F_q$. A forbidden triple has the form $f,rf,r^2f$ with nonconstant $r\in\mathbb F_q(x)$ and all three terms monic polynomials. In factorization coordinates the three vectors of valuations form an arithmetic progression with a nonzero integral step, which may have either sign in each coordinate. Write $\mathbb P_n$ for the uniform measure on the $q^n$ monic polynomials of degree $n$. The limiting proportion in balls exists with value $b$ exactly when $\mathbb P_n(E)\to b$: subtract successive ball counts for the forward implication and sum the geometric weights for the converse. All results also hold for all nonzero polynomials, since a fixed-scalar class gives an upper bound by the same averaging argument and our constructions may include all scalar multiples.

For $T\ge1$ let $S_T$ consist of the monic irreducibles of degree at most $T$, let $\|h\|=q^{\deg h}$, and set
\[
 c_T=\prod_{P\in S_T}(1-\|P\|^{-1}),\qquad
 v_T=\prod_{P\notin S_T}(1-\|P\|^{-2}).
\]
For subsets $A$ of the monic $S_T$-smooth polynomials containing no forbidden triple define
\[
 B_T=c_T\max_A\sum_{h\in A}\|h\|^{-1}.
\]
The maximum exists. Indeed, identify subsets with their indicator vectors in a countable product of two-point sets. Admissibility is closed, and the objective is continuous by convergence of the full Euler product $1/c_T$.

\paragraph{The extremal theorem.}
If $b$ is the supremum of natural densities of admissible sets, then
\[
 v_T B_T\le b\le B_T,\qquad
 0\le B_T-b\le q^{-T}/(q-1).
\]
In particular $b$ is computable, and there exists an admissible set of natural density exactly $b$.

\paragraph{Upper bound.}
For an admissible set $E$ of density $d$, write $f=gh$ with $g$ $S_T$-rough and $h$ $S_T$-smooth. The fiber $A_g=\{h:gh\in E\}$ is admissible. Average exact-layer proportions for $0\le n<N$. Their sum is
\[
 \sum_{\substack{g\ {\rm rough}\\\deg g<N}}\|g\|^{-1}
 \sum_{\substack{h\in A_g\\\deg h<N-\deg g}}\|h\|^{-1}
 \le (B_T/c_T)\sum_{\substack{g\ {\rm rough}\\\deg g<N}}\|g\|^{-1}.
\]
By finite inclusion--exclusion the rough proportion is exactly $c_T$ in every sufficiently high degree. Division by $N$ and passage to the limit give $d\le B_T$. Applying the sum over all scalar classes proves the same bound when their aggregate has density, whether or not each class individually has density.

\paragraph{Lower bound and computability.}
Use a maximizing $A$ uniformly in every rough fiber for which $g$ is squarefree and omit every other fiber; call this set $E_T$. In an allowed progression the three valuations outside $S_T$ would lie in $\{0,1\}$, so they agree coordinatewise. The common fiber excludes the progression. The density of $E_T$ is $e_T=v_T B_T$. To verify this, first impose a degree bound on $h$ and check squarefreeness only at primes up to a second cutoff. Each resulting condition is a finite Boolean combination of divisor events, and $\mathbb P_n(a\mid f)=\|a\|^{-1}$ once $n\ge\deg a$. The error from large smooth factors is bounded by $\sum_{\deg h>L,\, h\ {\rm smooth}}\|h\|^{-1}$ and that from large squares by $\sum_{\deg P>K}\|P\|^{-2}$. Both bounds hold uniformly in $n$ since the probability of divisibility is always at most the reciprocal norm. The product formula follows on letting the cutoffs grow. Since $B_T\le1$,
\[
 1-v_T\le\sum_{d>T}q^d q^{-2d}=q^{-T}/(q-1).
\]
Finite weighted independent-set optimization on $\deg h\le L$ bounds the unrestricted maximum from below; the preceding smooth tail bounds the deficit. For example the tail is at most $\sum_{n>L}(n+1)^{|S_T|}q^{-n}$. These bounds give an algorithm with any prescribed absolute accuracy.

\paragraph{A uniform estimate for the largest irreducible factor.}
Write $D(f)$ for the largest degree of an irreducible divisor ($D(1)=0$), and $F_m(n)=\mathbb P_n(D(f)\le m)$. For $m\ge1$ we have
\[
 |F_m(n)-F_m(n-1)|\le 1/m\quad(n\ge1).
\]
In fact, let $C(z)=(1-z)\sum_{n\ge0}F_m(n)z^n$. The Euler product over monic polynomials gives
\[
 C(z)=\prod_{\deg P>m}(1-(z/q)^{\deg P})
     =\exp\!\left(-\sum_{k>m}a_kz^k/k\right),\quad
 a_k=q^{-k}\!\sum_{\substack{d\mid k\\d>m}} dI_q(d)\in[0,1],
\]
where $I_q(d)$ counts degree-$d$ monic irreducibles and $\sum_{d\mid k}dI_q(d)=q^k$. If $c_n=[z^n]C(z)$, then $c_n=0$ for $1\le n\le m$. For $n>m$ logarithmic differentiation and coefficientwise majorization by the exponential with positive signs give
\[
 n|c_n|\le\sum_{j=0}^{n-1}|c_j|
 \le\exp\!\left(\sum_{k=m+1}^{n}1/k\right)\le n/m .
\]
The coefficient $c_n$ is precisely the first difference in question. Consequently $|F_m(n-k)-F_m(n)|\le k/m$ for $0\le k\le n$.

\paragraph{Pasting by largest prime degree.}
For each $j\ge1$, take the product set $E_j$ constructed at $T=j$, of density $e_j\to b$, and put $\epsilon_j=2^{-j}$. By the two uniform truncations above, there is a cylinder event $C_j$ with $\mathbb P_n(E_j\mathbin{\triangle} C_j)\le\epsilon_j$ for every $n$. Its indicator has the form $\sum_a u_a\mathbf1_{a\mid f}$ for finitely many monic $a$ by inclusion--exclusion (intersections use least common multiples). Write $s_j=\sum_a u_a/\|a\|$; then $|s_j-e_j|\le\epsilon_j$. Choose increasing positive integers $m_j>j$ exceeding every such $\deg a$ for $C_j$ and satisfying $\sum_a |u_a|\deg(a)/\|a\|\le m_j\epsilon_j$.

Accept $f$ by rule $E_j$ when $m_j<D(f)\le m_{j+1}$, and omit $D(f)\le m_1$. This single set avoids every forbidden triple. Indeed, at a prime of the greatest degree occurring anywhere in a putative triple all three valuations are zero or one: wherever present it is a largest prime and the rule there enforces squarefreeness. As the valuations are in arithmetic progression and one equals one, all equal one. Thus all three polynomials have the same largest degree, follow the same rule, and cannot form the triple.

For $m,n\ge m_j$, uniformity follows from exact division by $a$ and the shift estimate:
\[
 \mathbb P_n(C_j, D\le m)=
 \sum_a (u_a/\|a\|)F_m(n-\deg a)
 =s_jF_m(n)+O(\epsilon_j).
\]
The implicit constant is one. Replacing $C_j$ by $E_j$ and $s_j$ by $e_j$ increases it to three. The joint probability on the $j$th bin therefore differs from $e_j\mathbb P_n(m_j<D\le m_{j+1})$ by at most $6\epsilon_j$; for $n<m_j$ the bin is empty. For each fixed $m$, $F_m(n)\to0$, since the finitely many permitted prime factors give at most $(n+1)^{|S_m|}$ possible polynomials. Hence the total mass of any fixed initial segment of bins tends to zero. The sum of the remaining comparison errors is bounded by $6\sum_{j\ge J}\epsilon_j$, and $\sup_{j\ge J}|e_j-b|\to0$. Letting $n$ and then $J$ tend to infinity proves that the exact-layer proportions tend to $b$.

\paragraph{Effective realization and the greedy comparison.}
The realizing set can be chosen decidable. At stage $j$ use a finite optimizer whose omitted smooth tail, after multiplication by $c_j$, is less than $2^{-j}$. Deterministic tie breaking makes each optimizer computable. The cylinder truncations and the endpoints $m_j$ can then be chosen by the displayed effective bounds. This changes $e_j$ by at most $2^{-j}$ and hence preserves its limit. Factoring a given polynomial and finding its bin decides membership.

Let $G$ be the usual greedy set, consisting of polynomials whose valuations have only $0$ and $1$ as ternary digits. Its full density is the product of its one-prime masses. There is a natural-density improvement over $G$ if and only if for some $T$ there is a \emph{finite} admissible set of $S_T$-smooth polynomials whose reciprocal-norm sum strictly exceeds that of the entire local greedy set. Indeed, a strict local gain can be tensored with the greedy rule at every outside prime. In a progression all outside coordinates must then be constant; convergence of the outside product and its density follow as before, since any excluded exponent is at least two. Conversely, if $b>d(G)$, then for all sufficiently large $T$, $B_T$ exceeds the local greedy mass after normalization by $c_T$, because the omitted greedy product tends to one. Truncating a maximizer in degree preserves a strict inequality for all sufficiently large truncation degrees. Thus strict improvement admits a finite certificate, and exhaustive search will find one whenever it exists. The argument leaves the outcome of this comparison open.

\paragraph{Uniform means on long intervals of degrees.}
The same number $b$ is the largest upper Banach mean of the sequence of exact-degree proportions. In detail, for any admissible $E$ and fixed $T$ one has
\[
 \sum_{n=M}^{M+L-1}\mathbb P_n(E)\le LB_T+O_{q,T}(1)
 \qquad(M\ge0,\ L\ge1).
\]
For rough cores having degrees in $[M,M+L)$ the fiber bound and the eventually exact rough proportions give $LB_T+O_{q,T}(1)$. The contribution of earlier rough cores is at most
\[
 \sum_{s\ge1}\sum_{\substack{h\ {\rm smooth}\\\deg h\ge s}}\|h\|^{-1}
 =\sum_{h\ {\rm smooth}} (\deg h)\|h\|^{-1}<\infty ,
\]
by summing over $s=M-\deg g$ and bounding the rough shell proportion by one. Divide by $L$, take the supremum over the starting position $M$, then let $L$ and $T$ tend to infinity. The reverse extremal inequality follows from the constructed set with converging degree proportions.

\paragraph{A finite binary certificate.}
For $q=2$ choose $S_2$ comprising the two linear primes and the unique quadratic prime. On exponent triples $a+b+2c\le12$ the maximum of $\sum 2^{12-a-b-2c}$ is $14853$, attained by the truncated greedy set. The script \texttt{natural\_density\_check.py} exhausts the independent sets by branching on the earliest remaining vector. Including it deletes the third vertex of any progression whose first two in the chosen ordering have now been included. A branch can be pruned whenever its score plus the total available weight cannot improve the incumbent. All scores are integers. Assigning full mass to the untested smooth tail gives
\[
 b_2\le 1-\frac3{16}\left(\sum_{a+b+2c\le12}2^{-a-b-2c}
                   -\frac{14853}{4096}\right)
 =\frac{2797}{4096}=0.682861328125.
\]
The global greedy density is approximately $0.6483610761$ by the Euler product of Asada et al.

\section{A strict gap between natural density and upper density}
\label{sec:gap}
Let $b_q$ denote the supremum for sets whose degree-ball density converges and $u_q$ the supremum when a limsup is allowed. Ratios range over nonconstant rational functions, including those of degree zero. Proportions for monic polynomials suffice as in the preceding section.

\begin{theorem}\label{th:gap}
For every prime power $q$, $u_q>b_q$. For each $T\ge3$,
\[
 u_q-b_q\ \ge\ (1-q^{-1}) c_T q^{-3}
            -\frac{q^{-T}}{(T+1)(q-1)} .
\]
In particular $u_2-b_2>0.003660$. As $q\to\infty$ through prime powers, $\liminf q^3(u_q-b_q)\ge e^{-11/6}$.
\end{theorem}

\paragraph{Separate optimization at each cutoff.}
Let $J_T(t)$ be the largest admissible subset of the monic $S_T$-smooth polynomials of degree at most $t$. Define
\[
 U_T=(1-q^{-1})c_T\sum_{t\ge0}q^{-t}J_T(t).
\]
The upper-density inequalities $v_T U_T\le u_q\le U_T$ occur in the incoming report \emph{A counterexample to the geometric-progression density conjecture over finite fields}, section on computing the extremal upper density. Here is a direct proof for signed valuation steps. A ball through degree $N$ has $(q^{N+1}-1)/(q-1)$ monic elements. The number of rough cores of degree $N-t$ divided by $q^{N-t}$ tends to $c_T$. There are at most $J_T(t)$ admissible smooth factors per core. Dominated convergence gives the upper bound since $J_T(t)\le(t+1)^{|S_T|}$. For the lower bound use only squarefree rough cores; their degree proportions tend to $c_Tv_T$ by finite divisibility and the convergent square tail. Fix $N_j=2^{j+1}$, with an initial cutoff $N_0=0$. For each such core with degree in $(N_{j-1},N_j]$, permanently choose a maximizing smooth fiber at cutoff $N_j-\deg g$. Distinct cores cannot occur in a progression because their outside valuations are all in $\{0,1\}$. At $N_j$ the entire contribution from $t<N_j-N_{j-1}$ attains the finite-fiber bound. The summable tail tends to zero. This proves the lower estimate. Taking all possible smooth factors in the same count gives $U_T\le1$.

For any fixed admissible smooth set $A$, let $A(t)$ count its elements through degree $t$. A geometric series gives
\[
 \sum_{h\in A}\|h\|^{-1}=(1-q^{-1})\sum_{t\ge0}q^{-t}A(t).
\]
Each difference $J_T(t)-A(t)$ is a nonnegative integer. We will show that one is positive for $1\le t\le3$. Because $T\ge3$, the smooth elements through these cutoffs are all the monic polynomials there.

\paragraph{Why the three cutoffs are incompatible.}
Write $\mathcal L$ for the $q$ monic linears. Every monic rational-ratio progression has the form
\[
                  a d^2,\quad a d e,\quad a e^2,
\]
for distinct relatively prime monic $d,e$ and monic $a$. This follows by expressing the ratio in lowest terms. At cutoff three, $d,e\in\{1\}\cup\mathcal L$ and $\deg a\le1$. At cutoff two we must have $a=1$. For each pair of distinct diagonal entries $d^2,e^2$, its unique middle entry $de$ is required to be absent if both diagonals are retained. Retaining $k$ of $q+1$ diagonals costs at least
\[
                  q+1-k+\binom{k}{2}\ \ge\ q
\]
deletions, with distinct middle entries for distinct pairs. The bound is attained by deleting all $q$ linear squares. At cutoff one all $q+1$ entries are admissible. Thus any set optimizing both cutoffs one and two must delete precisely those $q$ squares. For each $a,d\in\mathcal L$ with $a\ne d$, the lower-degree entries $a,ad$ have been retained. Their progression forces deletion of $ad^2$. These $q(q-1)$ distinct cubics give at least $q^2$ deletions through degree three.

There is a set at cutoff three with only $q^2-1$ deletions. Start with the deletion list consisting of every $P^2$ and every $ad^2$ for distinct $a,d\in\mathcal L$. Choose distinct $X,Y\in\mathcal L$, restore $Y^2$ and $YX^2$, and delete $Y$. To check it, group the triples by their common factor $a$ in the displayed parametrization. For $a=1$ the retained diagonals are $1,Y^2$, whose middle is the missing $Y$. For $a=Y$ the retained diagonals are $Y^3,YX^2$, whose middle is the missing $XY^2$. For each remaining linear $a$, the only retained diagonals are $a,a^3$, whose middle $a^2$ is missing. This exhausts the possibilities, including ratios of degree zero.

\paragraph{The gap in the limit.}
The integer loss has weight at least $q^{-3}$ in the series. Maximizing over $A$ gives $U_T-B_T\ge(1-q^{-1})c_Tq^{-3}$. Together with $u_q\ge v_T U_T$ and $b_q\le B_T$ this proves the bound, since $I_q(d)\le q^d/d$ gives
\[
 1-v_T\le\sum_{d>T} I_q(d)q^{-2d}
       \le q^{-T}/((T+1)(q-1)).
\]
The product defining $c_T$ decays sufficiently slowly: its first factor is $(1-1/q)^q\ge1/4$, and for $d\ge2$ Bernoulli's inequality gives $(1-q^{-d})^{I_q(d)}\ge1-1/d$. Hence $c_T\ge1/(4T)$, and sufficiently large $T$ gives a positive gap for every $q$. For the asymptotic, take $T=3$ and use $I_q(1)=q$, $I_q(2)=(q^2-q)/2$, $I_q(3)=(q^3-q)/3$. Then $c_3\to e^{-11/6}$.\hfill$\square$

\begin{center}\small
\begin{tabular}{c r l@{\qquad}c r l}\toprule
$q$&$T$&Certified gap exceeds&$q$&$T$&Certified gap exceeds\\\midrule
2&8&0.003660&5&4&0.000715\\
3&5&0.002149&7&4&0.000297\\
4&4&0.001192&&&\\\bottomrule
\end{tabular}
\end{center}
The script \texttt{check\_gap.py} verifies the table with rational arithmetic. It also exhausts the $2^{15}$ subsets in the binary degree-three ball: its optimal sizes through degrees $0,1,2,3$ are $1,3,5,12$. Preserving the optima through degree two allows at most $11$ at degree three.

\paragraph{Polynomial ratios.}
The same separation holds if $r$ must be a nonconstant polynomial. Use admissibility for this ratio convention in both optimizations; outside-squarefreeness still fixes the rough core. At cutoff two the unique monic constant $1$ meets every forbidden triple, so one deletion suffices and is necessary. A set optimal at cutoff one retains $1$ and all $q$ linears and must delete all $q$ linear squares at cutoff two. This proves incompatibility already through degree two and replaces $q^{-3}$ in Theorem~\ref{th:gap} by $q^{-2}$ for this convention.

\section{A quadratic Jensen theorem for a nonabelian Artin factor}
\label{sec:eta-jensen}
Let $F=\mathbb Q(\alpha)$ with $\alpha^3-\alpha-1=0$. The polynomial is
irreducible over $\mathbb Q$ and has discriminant $-23$. Its splitting
field has group $S_3$. Write $\rho$ for the two-dimensional summand of its
three-letter permutation representation. The classical weight-one identity is
\[
 L(s,\rho)=\frac{\zeta_F(s)}{\zeta(s)}=L(s,f),
 \qquad f(z)=\eta(z)\eta(23z).
\]
Here $L(s,f)$ uses the Fourier coefficients of the normalized cusp form.
See van der Blij's result in D. Zagier, \emph{Elliptic Modular Forms and
Their Applications}, Proposition 14 and the discussion following it,
\href{https://www.its.caltech.edu/~matilde/Zagier123ModularForms.pdf}{pp.~43--44}.
In particular the trace at an unramified prime is the number of roots of
$X^3-X-1$ in $\mathbb F_p$ minus one.
Put
\[
 \Lambda_f(s)=\left(\frac{\sqrt{23}}{2\pi}\right)^s\Gamma(s)L(s,f),
 \qquad
 \Lambda_f(1/2+z)=\sum_{n\ge0}\frac{\gamma_n}{n!}z^{2n}.
\]
\textbf{Theorem.} For every $n\ge0$, the polynomial
$\gamma_n+2\gamma_{n+1}X+\gamma_{n+2}X^2$ has two distinct negative
real roots. The conclusion holds independently of Dirichlet GRH.

\textit{Proof.}
After the change of variables $y=e^u/\sqrt{23}$ in the Mellin integral
of $f(iy)$, we obtain
\[
 \Lambda_f(1/2+z)=\int_{\mathbb R}e^{zu}K(u)\,du,
 \qquad K(u)=e^{u/2}\eta(ie^u/\sqrt{23})\eta(i\sqrt{23}e^u).
\]
The integral agrees with the completed Dirichlet series in its initial
half-plane, then gives its entire continuation by the decay below.
Write $h(v)=v/4+\log\eta(ie^v)$ and $a=\frac12\log23$. Then
$\log K(u)=h(u-a)+h(u+a)$. Eta inversion shows that $h$ is even.
For $v\ge0$, the product for $\eta$ gives, with $Y=2\pi e^v$,
\[
 h(v)=\frac v4-\frac Y{24}
       -\sum_{m\ge1}\sigma_{-1}(m)e^{-mY},\qquad
 \sigma_{-1}(m)=\sum_{d\mid m}d^{-1}.
\]
The series and its first four derivatives converge locally uniformly.
The fourth derivative in $v$ of $e^{-mY}$ equals
$e^{-q}q((q-6)q^2+7q-1)$ with $q=mY\ge 2\pi>6$. Consequently
$h^{(4)}(v)<0$ for all real $v$, using evenness for $v<0$.

More generally, let $g(u)=\sum_j r_j[h(u-a_j)+h(u+a_j)]$ for any
nonempty finite list of $a_j\in\mathbb R$, $r_j>0$, and put $K=e^g$.
Then $g$ is even and $g^{(4)}<0$. Thus $g'''(u)<0$ for $u>0$ and
$(u g''-g')'=u g'''<0$. Since $g'(0)=0$,
\[
 \left(\log K(\sqrt{x})\right)''
 =\frac{u g''(u)-g'(u)}{4u^3}<0,\qquad u=\sqrt x>0 .
\]
At both infinities $K$ decays at least as fast as $\exp(-c e^{|u|})$
for some $c>0$, by the same product and symmetry. The preceding
calculation follows the eta-kernel curvature argument of J. B. Conrey
and A. Ghosh, \emph{Tur\'an inequalities and zeros of Dirichlet series
associated with certain cusp forms}, Trans. AMS 342 (1994), 407--419,
\href{https://aimath.org/~kaur/publications/29.pdf}{Secs.~4--7}.
The symmetry of the translations preserves the fourth-derivative sign
for unequal eta dilations.

For completeness, write $V(x)=-\log K(\sqrt{x})$ and
$M(p)=\int_0^\infty x^{p-1}e^{-V(x)}\,dx$ for $p>0$. Under probability
density $x^{p-1}e^{-V(x)}/M(p)$ denote expectation by $E_p$.
Integration by parts gives
$E_{p+1}[V']=p M(p)/M(p+1)$. Since $V'$ is strictly increasing,
\[
 E_{p+2}[V']-E_{p+1}[V']
 =\frac{\operatorname{Cov}_{p+1}(x,V'(x))}{E_{p+1}[x]}>0.
\]
The last sign follows by using two independent draws and averaging
$(x-y)(V'(x)-V'(y))/2$. Hence the consecutive ratios of $M(p)/\Gamma(p)$
strictly decrease as $p$ increases by one. Since
\[
 \gamma_n=\frac{n!}{(2n)!}M(n+1/2)
 =\frac{\sqrt\pi}{4^n}\frac{M(n+1/2)}{\Gamma(n+1/2)},
\]
we have $\gamma_{n+1}^2>\gamma_n\gamma_{n+2}$. All coefficients are
positive. This proves the quadratic assertion.\hfill$\square$

The same calculation applies to every positive eta product with
exponents paired under $d\leftrightarrow N/d$, on centering its
Mellin kernel at weight divided by two. In general this kernel need not
come from an eigenform. For the conductor-$23$ form, the open critical-line
statement would imply real-rootedness in every Jensen degree:
$w\mapsto\Lambda_f(1/2+\sqrt w)$ has order at most $1/2$, and critical-line
zeros correspond to negative real $w$. Hadamard factorization and Jensen's
criterion then give this implication and its converse. The theorem proves
the degree-two part for all shifts. The kernel argument supplements the
prime-sum reduction. It makes no use of its arithmetic estimate.

\subsection{A counterexample to the general cubic kernel criterion}
A possible next step would apply Theorem~1 of D. K. Dimitrov and
F. R. Lucas, \emph{Higher order Tur\'an inequalities for the Riemann
$\xi$-function}, Proc. AMS 139 (2011), 1013--1022,
\href{https://doi.org/10.1090/S0002-9939-2010-10515-4}{p.~1017}.
Together with their normalization and Lemma~1, the printed theorem implies
that every cubic Jensen polynomial is hyperbolic for an admissible
kernel with $(\log K(\sqrt{x}))''<0$ for $x>0$. The following counterexample
addresses this general assertion. It leaves the individual questions for
the Riemann and conductor-$23$ kernels separate.

\textbf{Counterexample.} Put $K_\delta(u)=\exp(-u^2-\delta u^6)$ with
$\delta=10^{-8}$. Its first cubic Jensen polynomial has a nonreal
conjugate pair. This kernel is positive, even and entire, decreases
strictly for $u>0$, and every derivative is a polynomial times $K_\delta$,
hence $O(e^{-|u|^3})$ at infinity. Moreover
$(\log K_\delta(\sqrt{x}))''=-6\delta x<0$.

\textit{Proof.} Use the coefficients $\gamma_n$ above with this new
kernel and put $c_n=4^n\gamma_n/\sqrt\pi$. If $p=n+1/2$ and
$X$ has gamma density $x^{p-1}e^{-x}/\Gamma(p)$, then
$c_n=\mathbb E e^{-\delta X^3}$. With $(p)_m=p(p+1)\cdots(p+m-1)$,
Taylor's inequalities for $e^{-y}$ on $y\ge0$ give
\[
 1-\delta(p)_3\le c_n\le1-\delta(p)_3+
                   \tfrac12\delta^2(p)_6 .
\]
The rescaled Jensen polynomial is
$P(Z)=c_0+3c_1Z+3c_2 Z^2+c_3 Z^3$. These rescalings preserve
real-rootedness. Since $(n+1/2)_3$ is monic cubic in $n$,
its third forward difference is $6$, and $P(-1)=6\delta+O(\delta^2)$.
In $P(W-1)$ the respective coefficients of $W^3,W^2,W,1$ are
$1+O(\delta),O(\delta),O(\delta),6\delta+O(\delta^2)$. The discriminant
is therefore $-972\delta^2+O(\delta^3)$ as $\delta\downarrow0$.
For the stated rational value the accompanying checker uses the displayed
rational endpoints to certify
\[
                 -975<\operatorname{disc}(P)/\delta^2<-969.
\]
Explicitly it encloses $T_1=c_1^2-c_0c_2$, $T_2=c_2^2-c_1c_3$, and
$U=c_1c_2-c_0c_3$ by subtraction of endpoint products. All three lower
endpoints are positive. The final enclosure follows from
$\operatorname{disc}(P)=27(4T_1T_2-U^2)$. \hfill$\square$

One can locate a difficulty in the mean-value argument immediately
before equation (14) on pp.~1020--1021. There
$A(y,z)=(z^2-y^2)K(y)K(z)[B(y)-B(z)]$ with
$B(y)=K'(y)/(yK(y))$. For $y>0$ and $B'(y)<0$, $A(y,z)$ vanishes
quadratically as $z\to y$. Hence $\partial_y A(y,z)/A(y,z)$ has a
simple pole at $z=y$. Its derivative on either open side cannot justify
the mean-value step for an interval with $z_1<y<z_2$.
An argument for the cubic inequalities of a particular arithmetic
kernel must use additional information. The counterexample is a
Mellin-moment calculation and makes no assertion about zeros of a
particular number-theoretic $L$-function.

The obstruction persists with globally strict fourth-derivative curvature.
Take $\widehat K_\delta(u)=\exp[-u^2-\delta(u^4+u^6)]$. Its log has fourth
derivative $-24\delta-360\delta u^2<0$, as does the log of the symmetric
eta kernel above. In the gamma expectation the perturbing polynomial is now
$X^2+X^3$, so the first-order correction to $c_n$ is
$-\delta[(p)_2+(p)_3]$. Its third forward difference is again $-6\delta$.
For every fixed shift $n$ the same translated-cubic calculation gives
discriminant $-972\delta^2+O_n(\delta^3)$ in the $c_n$ normalization.
A single sufficiently small $\delta>0$ therefore violates the cubic
inequalities at any prescribed finite set of shifts. This distinguishes
the information provided by the eta curvature estimate from what a
degree-three argument requires.

\subsection{A necessary condition on perturbations of the Gaussian}
The calculation also classifies the possible first-order changes of the kernel if all the cubic tests are to survive near a Gaussian. Let $V:[0,\infty)\to\mathbb R$ be continuous, bounded below and of at most polynomial growth. For $\epsilon>0$ use
\[
 K_\epsilon(u)=\exp\{-u^2-\epsilon V(u^2)\}.
\]
Suppose there is a sequence $\epsilon_i\downarrow0$ along which the cubic Jensen polynomials at every shift are real-rooted. Then $V(x)$ is a polynomial of degree at most two. The same conclusion holds if for each fixed shift the sequence may be chosen separately. The statement is a necessary condition on this family of perturbations.

Indeed, for $X_n$ of gamma law with shape $n+1/2$ and rate one, put $m_n=\mathbb E V(X_n)$. With the rescaling of the preceding proof,
$c_n(\epsilon)=1-\epsilon m_n+O_n(\epsilon^2)$. The assumptions bound the integrated quadratic Taylor remainder: for $V\ge0$ use $\epsilon^2 V^2/2$, and for $V<0$ multiply this bound by $\exp(\epsilon|\inf V|)$. At shift $n$ the translated cubic has constant term $\epsilon\Delta^3m_n+O_n(\epsilon^2)$, where $\Delta m_n=m_{n+1}-m_n$. Its quadratic and linear coefficients are $O_n(\epsilon)$ and its leading coefficient is $1+O_n(\epsilon)$. Its discriminant begins with
\[
                   -27\epsilon^2(\Delta^3m_n)^2+O_n(\epsilon^3).
\]
Thus $\Delta^3m_n=0$ at each $n$. The sequence $m_n$ is quadratic in $n$. There is a unique polynomial $Q$ of degree at most two with $\mathbb E Q(X_n)=m_n$ for all $n$, since the expectations of $1,x,x^2$ are $1,p,p(p+1)$ at $p=n+1/2$.

Now $(V-Q)x^{-1/2}e^{-x}\,dx$ on $(0,\infty)$ is a finite signed measure with all moments zero. Its Laplace transform, in the convention $e^{zx}$, is analytic for $\Re z<1$ by the growth assumption and all its derivatives at zero vanish. Analytic continuation followed by uniqueness of Fourier transforms of finite measures gives the zero measure. Continuity gives $V=Q$ everywhere. For any other $V$ in this class there is at least one fixed shift whose cubic has a nonreal pair for every sufficiently small positive perturbation. This explains the third finite difference in the explicit sextic example.

\subsection{An independent proof through degree sixteen, at every shift}
There is another route for this particular Artin factor that uses the
location of its low zeros. Define $G(w)=\Lambda_f(1/2+\sqrt w)$, with
analyticity at zero given by evenness. We prove
\[
 J^{d,n}(X)=\sum_{j=0}^d \binom{d}{j}\gamma_{n+j}X^j
 \quad\text{has only negative real zeros for } 1\le d\le16,\quad n\ge0 .
\]
There are two ingredients.

\textit{A zero-free rectangle.}
Put $a=\Re(s)-1/2$, $t=\Im(s)$. For the conductor-$23$ eta kernel of Section~\ref{sec:eta-jensen},
\[
 \tfrac12\Re\Lambda_f(s)=\int_0^\infty K(u)\cosh(au)\cos(tu)\,du.
\]
On $|a|\le1/2$ and $|t|\le4$, the accompanying \texttt{check\_height.py}
proves the right side exceeds $1/1000$ using integer bounds.
Here are the steps of the certificate. Subdivide $0\le u\le4$ in steps
of $1/1024$ and $0\le t\le4$ in steps of $1/256$. In each $u$ interval
bound $q=\exp[-(2\pi/\sqrt{23})e^u]$ using
$1.3101<2\pi/\sqrt{23}<1.3102$. In
$K=e^{u/2}q\prod_{m\ge1}(1-q^m)(1-q^{23m})$ keep the first eight
factors of the first product. The product of the remaining factors is
between $1-2q^9/(1-q)$ and $1$, by
$\prod(1-v_i)\ge1-\sum v_i$ for $0\le v_i\le1$. Bound the phase cosine
from below by evaluating its degree-$62$ lower Taylor polynomial at
$\lfloor512 t_{\min}u_{\min}\rfloor/512$ and subtracting the maximum
phase distance from this point in the cell. This uses $|\cos'|\le1$.
A positive cosine lower bound is multiplied by the lower bound for
$K$, a negative one by the upper bound for $K\cosh(u/2)$.
The sum of these interval contributions is a lower integral for every $a$
and every $t$ in that cell of heights. Beyond $u=4$ the absolute integrand
is bounded by $e^u e^{-e^u}$, giving tail $e^{-e^4}<2^{-64}$.
All exponentials in the certificate use a rational Taylor bound on
$[0,1]$ followed by seven squarings with outward integer rounding.
Thus no floating-point or library zero-finder is required.

\textit{The sector argument.}
The Euler product and functional equation put the zeros of the completed
Artin function in $|a|\le1/2$. The zeros of $G$ are $w=(a+it)^2$.
If all zeros with $|t|<T$ are on the line $a=0$, or merely satisfy
$|a|/|t|\le1/(2T)$, the zeros of $G$ lie in the closed sector about
the negative real ray with half-angle $\theta=2\arctan(1/(2T))$,
provided $T>1/2$ and $\Lambda_f(1/2)\ne0$. The growth estimate
$\log\max_{|z|\le R}|\Lambda_f(1/2+z)|=O(R\log(2+R))$ follows by
integrating $e^{R|u|-c e^{|u|}}$. Hence $G$ has order at most $1/2$ and
is the locally uniform limit of its genus-zero partial products. Group
conjugate zeros to obtain real polynomials. The closed sector is convex;
Gauss--Lucas and locally uniform convergence place the zeros of every
$G^{(n)}$ in the same sector. All $G^{(n)}$ have positive Taylor
coefficients by the Mellin formula.

The sector theorem of Obreschkoff and Chasse now gives real-rooted
Jensen polynomials for each $G^{(n)}$ through degree
\[
   d\le\csc^2\theta=\left(T+\frac{1}{4T}\right)^2.
\]
See M. Chasse, \emph{Laguerre multiplier sequences and sector properties
of entire functions}, Complex Var. Elliptic Equ. 58 (2013), 875--885,
\href{https://doi.org/10.1080/17476933.2011.584250}{journal version}, or
Theorem~171 and the application to derivatives in Theorem~173 of his
\href{https://people.kth.se/~chasse/dissertation.pdf}{dissertation}.
Taking $T=4$ gives $d\le\lfloor(65/16)^2\rfloor=16$.
Positivity of the coefficients places all real roots on the negative ray.

\textit{Larger ranges from external computations.}
The same sector calculation with $T=18.08$ gives $d\le327$ from the
published interval data. LMFDB identifies this function with classical
newform $23.1.b.a$, lists ten positive low ordinates from $5.1156833$
through $18.0806617$, and records rank zero. Its documentation for modular
forms describes a complete list of zero disks within the listed range.
The displayed radii give $|a|<10^{-3}$ and $|t|>5$ in those disks.
This suffices for $|a|/|t|<1/(2T)$ even without a proof that each disk's
zero lies exactly on the line. See the
\href{https://www.lmfdb.org/L/2/23/23.22/c0/0/0}{record} and
\href{https://www.lmfdb.org/L/2/23/23.22/c0/0/0/Reliability}{precision
and completeness statement}.

In his thesis Nathan Ng reports $1{,}736$ positive zeros and verification
of RH for this Artin factor up to $T=972.2$ using Tollis's zero-counting
theorem. Taking that computation as input gives $d\le945{,}173$. The
thesis illustrates the Gram-block test in detail at level $31$ and
tabulates its outcome for level $23$. We did not rerun the height-$972.2$
computation. See \emph{Limiting Distributions and Zeros of Artin
$L$-Functions}, \href{https://www.cs.uleth.ca/~nathanng/RESEARCH/phd.thesis.pdf}{Secs.~4.3.4--4.3.5}.
None of these finite ranges uses the cubic kernel assertion above.

\section{The least degree for odd projective special linear groups}
\label{sec:oddpsl}
The earlier \texttt{psl2\_grh\_transfer.zip} studies $\mathrm{PSL}_2$ in characteristic two. Its third further question asks for the least dimension of a nonnegative character containing every nontrivial irreducible in odd characteristic and for optimal subfield degrees. Nonnegativity here means nonnegativity as a function on group elements. We answer the first question for every odd prime power $q\ge5$ and the single-subfield spectral question for every prime $q\ge127$.

\paragraph{The least dimension.}
Let $G=\mathrm{PSL}_2(\mathbb F_q)$ and $S=\sum_{\chi\in\mathrm{Irr}(G)}\chi$, the classical Gelfand character. The unique minimizer is $S$. Set $D=q(q+1)/2$. The complete set of weights is
\[
\begin{array}{c|cccc}
q\bmod4 &\text{identity}&\text{nontrivial unipotent}
 &\text{nonidentity split }t&\text{nonidentity nonsplit }t\\ \hline
1&D+1&1&1+\delta(t)&1\\
3&D&0&1&1-\delta(t).
\end{array}
\]
Here $\delta$ denotes the quadratic character of whichever projective torus has even order. The split and nonsplit projective tori $T_s,T_n$ have orders $a=(q-1)/2,b=(q+1)/2$, respectively. Toral involutions are included.

\paragraph{Proof.}
Use the ordinary character table in B. B\"ohmler, N. Farrell and C. Lassueur, \emph{Trivial source character tables of $\mathrm{SL}_2(q)$}, Table~1, restricting to rows trivial on $-I$ (\href{https://doi.org/10.14760/OWP-2022-02}{Oberwolfach preprint}). Besides $1$ and the Steinberg character of degree $q$, the principal series is indexed by nontrivial characters $\alpha$ of $T_s$ up to inversion. An inverse pair contributes one degree-$q+1$ row, with values $1$, $\alpha(t)+\alpha(t)^{-1}$ and $0$ on the three nonidentity types above. For the cuspidal series use nontrivial characters $\beta$ of $T_n$ up to inversion. A pair contributes one row of degree $q-1$ with values $-1$, $0$ and $-\beta(t)-\beta(t)^{-1}$. A quadratic parameter contributes two exceptional rows of half the corresponding degree. Their sum has exactly the displayed series values for this parameter. Steinberg has values $0,1,-1$.

If $q=4k+1$ the numbers of ordinary inverse pairs in the two series are $k-1,k$, and the exceptional pair is principal. If $q=4k+3$ the numbers are $k,k$, and the exceptional pair is cuspidal. Summing the degrees and the unipotent values gives the first two columns. For nontrivial $t$ the sum of all nontrivial characters of its cyclic torus is $-1$. A quadratic parameter, when present, contributes twice its value to the parameter sum in the character table, hence one extra copy beyond this $-1$. Adding the trivial and Steinberg terms gives the remaining columns. In particular $S(g)\ge0$ everywhere. Every nonzero, pointwise-nonnegative genuine character has strictly positive mean, since its value at the identity is its dimension. Its trivial multiplicity is therefore at least one. Covering all nontrivial irreducibles then forces degree at least $S(1)$, with equality exactly for $S$.

\paragraph{A Hecke realization when $q\equiv1\pmod4$.}
Let $N=T_s\rtimes\langle w\rangle$ be the dihedral normalizer, with $w$ acting by inversion. Extend $\delta$ to $\lambda$ on $N$ by $\lambda(w)=1$. Then
\[
 S=1+\operatorname{Ind}_N^G\lambda .
\]
Indeed the induced character has degree $D$ and vanishes on unipotents and nonsplit elements. All reflections in $N$ are involutions, which are split in this congruence class. Their total $\lambda$-sum is zero. At a split noninvolution the class intersects $T_s$ in $t,t^{-1}$ and its centralizer has order $a$. At the split involution the intersection with $T_s$ is a singleton and its centralizer has order $2a$. In both cases the induction formula gives $\delta(t)$. Thus for a $G$-extension $M/\mathbb Q$, $L(s,S)$ is the product of $\zeta_{\mathbb Q}(s)$ and a nontrivial quadratic Hecke $L$-function over $M^N$. This gives unconditional holomorphy in the open critical strip and a concrete nonabelian Frobenius weight for the one-sided analytic approach. A critical-line theorem for that Hecke function would still require new arithmetic estimates.

\paragraph{The smallest spectrally detecting single subfield.}
A subgroup $H$ is \emph{rich} if $\operatorname{Ind}_H^G1$ contains every irreducible, equivalently every irreducible has an $H$-fixed vector. In the terminology of the prior report, $M^H$ is then spectrally detecting. For every prime $q\ge127$,
\[
 \min_{H\ \mathrm{rich}} [G:H]=q(q-1).
\]
The nonsplit torus attains this minimum. To see that it is rich for every odd $q\ge5$, average each row of the ordinary table over $T_n$. The trivial, Steinberg, and each exceptional row have exactly one fixed dimension (on the nonsplit torus a cuspidal exceptional row has values $-\delta(t)$ away from the identity). Each ordinary principal or cuspidal row has two fixed dimensions. The latter computation uses $\sum_{t\ne1}\beta(t)=-1$.

Conversely, the dimension bound forces $|H|\le q-1$, since $|G|=(q-1)D$. Equality is impossible for $q\equiv1\pmod4$ since the least dimension is $D+1$. For $q\equiv3\pmod4$ equality would imply $\operatorname{Ind}_H^G1=S$. At a generator $t$ of $T_n$, $S(t)=2$, so $H$ would contain a conjugate of $t$. But $b=(q+1)/2$ fails to divide $q-1=2b-2$ for $b>2$. Thus $|H|<q-1$ in either case.

For prime $q\ge127$, Dickson's subgroup classification (Guralnick and Zieve, \emph{Polynomials with PSL(2) monodromy}, Appendix~A, \href{https://annals.math.princeton.edu/wp-content/uploads/annals-v172-n2-p11-p.pdf}{Theorem~A.1}) now leaves few possibilities. Since $|H|<q$, it has order prime to the defining characteristic. There are no proper subfield groups at prime $q$. The cyclic possibilities have order at most $b$. A dihedral group of order $2d$ has $d\mid a$ or $d\mid b$. The full normalizers have order $q-1$ or $q+1$, already excluded. Every smaller dihedral subgroup has order at most $b$. The exceptional groups $A_4,S_4,A_5$ have order at most $60<b$. These bounds prove the minimum. The restriction on $q$ matters. For example $\mathrm{PSL}_2(109)$ has a rich $A_5$ subgroup of order $60>55$.

\paragraph{Relation to earlier work and to zeros.}
Rich subgroups, their degree-sum bound, and finite examples are studied by Breuer, H\'ethelyi, Horv\'ath and K\"ulshammer, \emph{On Frobenius graphs of diameter 3 for finite groups}, Corollary~2.8 and Tables~2--3 (\href{https://doi.org/10.1016/j.jalgebra.2024.11.014}{J. Algebra 666 (2025), 507--529}). The large-prime formula above answers the optimization posed by the incoming report. No historical priority is claimed. Spectral coverage concerns which factors occur in the Artin factorization. Under Artin holomorphy in the open strip it gives an equivalence between RH for $\zeta_{M^{T_n}}$ and RH for $\zeta_M$, because all local zero orders are then nonnegative. Without that input a meromorphic constituent may have a negative order. The finite-group optimality proof supplies no sign control on those orders and no bound on the smoothed sum over primes.

\paragraph{Independent finite check.}
\texttt{check\_odd\_psl2.py} verifies the full weight for $q=5,7,11,13,17,19$ by exact arithmetic on projective matrices, independently of the character table. Writing $S=\sum m_\chi\chi$ for the proposed class function and taking unnormalized convolution, character orthogonality gives $(S*S)(g)=|G|\sum m_\chi^2\chi(g)/\chi(1)$. The classical commutator-count formula has the same right hand side with every $m_\chi^2$ replaced by one: average each conjugate of an irreducible matrix over $G$, obtaining the scalar trace divided by degree, and apply character orthogonality. The checker counts commutators and verifies this convolution equality on every class. It checks positive definiteness of convolution by $S$ on class functions by rational elimination of the Gram matrix on conjugacy-class indicators. Its eigenvalues on irreducible characters are $|G|m_\chi/\chi(1)$, hence every $m_\chi=+1$. It also verifies the induced-character identity in each applicable example. The degree minima there for the class-function problem are respectively $16,28,66,92,154,190$.

\section{Subfield degrees over every odd finite field}
\label{sec:oddpowers}
With the notation of Section~\ref{sec:oddpsl}, one can determine the least degree of a spectrally detecting single subfield for every odd prime power $q\ge5$. Let $h(q)$ be the maximum order of a rich subgroup of $\mathrm{PSL}_2(q)$. Except at four small fields, the answer is
\[
 h(q)=
 \begin{cases}
 q-\sqrt q,&q=s^2,\ s\equiv3\pmod4,\\
 q-q^{3/4},&q\text{ is a fourth power},\\
 (q+1)/2,&\text{otherwise}.
 \end{cases}
 \qquad
 \min_{H\ {\rm rich}} [G:H]=\frac{q(q^2-1)}{2h(q)} .
\]
The four exceptions are $h(19)=12$, $h(41)=24$, $h(101)=h(109)=60$. The index formula always uses the actual value of $h(q)$. In the first two cases an affine subgroup attains equality. For example $h(121)=110$, $h(169)=85$, and $h(625)=500$.

\paragraph{Reduction to an additive hyperplane.}
Write $q=p^f$ and $b=(q+1)/2$. The previous section proves $b\le h(q)<q-1$ for every odd $q\ge5$. Dickson's list, in the form of Theorem~A.1 cited there, excludes all exceptional groups from the order interval $(b,q-1)$ when $q\ge121$, since then $b>60$. For smaller $q$ we test these groups below. A proper subgroup of a torus normalizer has order at most $b$. A subfield $\mathrm{PSL}_2(t)$ with $q=t^j$ has order $(t^3-t)/2$. For $j=2$ its order exceeds $q-1$, and for $j\ge3$ it is less than $b$. A subfield $\mathrm{PGL}_2(t)$ requires even $j$. Its order $t^3-t$ exceeds $q-1$ for $j=2$ and is less than $b$ for $j\ge4$.

For the affine types on Dickson's list, a nontrivial normal elementary abelian $p$-subgroup has a unique common fixed point on the projective line: commuting nonidentity unipotents fix the same point. Its normalizer preserves that point, so the whole group lies in a Borel. The prime-to-$p$ subgroups of the affine type are cyclic of order at most $(q-1)/2$. It therefore remains to consider subgroups in a Borel group $B=U\rtimes T_s$, where $U=(\mathbb F_q,+)$ and $T_s=(\mathbb F_q^\times)^2$ acts by multiplication. For $H\le B$ let $A=H\cap U$ and $C$ be the image in $T_s$. If $A=U$ then $|H|\ge q$, and if $A=0$ then $|H|\le(q-1)/2$. When $C=1$, a proper additive subgroup has order at most $q/p$. Otherwise lift a generator of $C$ to an affine map and translate its unique fixed point to zero. This realizes $H=A\rtimes C$. Put $\mathbb F_r=\mathbb F_p(C)$. The space $A$ is $\mathbb F_r$-linear, so $|A|=q/r^m$ for some $m\ge1$. The condition $|H|>b$ gives $|C|/r^m>1/2$, while $|C|\mid r-1$. It follows that $m=1$ and $C=\mathbb F_r^\times$. Thus $A$ is an $\mathbb F_r$-hyperplane. Since $C$ consists of squares in $\mathbb F_q^\times$, $k=[\mathbb F_q:\mathbb F_r]$ must be even. In particular $q\equiv1\pmod4$.

\paragraph{The exceptional constituent test.}
We prove that $A\rtimes\mathbb F_r^\times$ is rich precisely when
\[
              4\mid\frac{q-1}{r-1}.
\]
Identify the additive characters of $U$ with $\mathbb F_q$ via the absolute trace pairing and a nontrivial additive character of $\mathbb F_p$. The annihilator $A^\perp$ is an $\mathbb F_r$-line, by nondegeneracy of the trace pairing. Its nonzero elements form one free $C$-orbit, contained in one of the two $T_s$-orbits of nontrivial additive characters. Whenever a representation contains a weight in this free orbit, it has an $H$-invariant: take any nonzero vector in that weight space and sum its translates by $C$. The summands have distinct $U$-weights and are all $A$-fixed.

Every ordinary principal, ordinary cuspidal, and Steinberg row contains all nontrivial $U$-weights. Indeed their restrictions to $U$ are, respectively, $\mathrm{reg}_U+1$, $\mathrm{reg}_U-1$ and $\mathrm{reg}_U$, by the unipotent column of the ordinary table. The exceptional pair $E_+,E_-$ consists of the two constituents of $\operatorname{Ind}_B^G\delta$, where $\delta$ is the quadratic character of $T_s$ inflated to $B$. Each has a $\delta$-line on restriction to $B$ by Frobenius reciprocity. Their combined restriction to $U$ is $\mathrm{reg}_U+1$. Hence each has exactly one trivial $U$-line. Each one's other weights fill one of the two $T_s$-orbits, since these have size $(q-1)/2$ and the degree of each exceptional row is $(q+1)/2$. The representation meeting $A^\perp\setminus\{0\}$ has an invariant by the free-orbit argument. In the other, the $A$-invariants are exactly its $\delta$-line. This proves richness if and only if $\delta|_C=1$. The kernel of $\delta$ in $\mathbb F_q^\times$ is the subgroup of fourth powers, so this is exactly the stated divisibility condition.

\paragraph{Maximization.}
For each admissible $r$ the affine subgroup has order $q(1-1/r)$, so choose the largest such $r$, equivalently the smallest admissible even $k$ with $q=r^k$. When $r\equiv3\pmod4$, the sum $1+r+\cdots+r^{k-1}$ is divisible by four for every even $k$. When $r\equiv1\pmod4$, divisibility is equivalent to $4\mid k$. Thus $k=2$ works precisely in the first displayed case. If $4\mid f$ then $k=4$ works, and the square root of $q$ is itself an odd square. If $f$ is odd there is no even $k$. Finally, if $f\equiv2\pmod4$ and $k=2$ fails, then $p\equiv1\pmod4$ and every possible even $k\mid f$ is twice an odd number, so none works. Each of the admissible maximizing affine orders exceeds $b$ whenever the corresponding case occurs ($q\ge9$). This proves the displayed baseline, subject only to checking exceptional subgroups of larger order when $q<121$.

\paragraph{The finite exceptional check.}
Write $h_0$ for the baseline in the cases display. For each odd prime power $5\le q<121$, retain only $K\in\{A_4,S_4,A_5\}$ on Dickson's list with $h_0(q)<|K|<q-1$. For odd $q$, $A_4$ always occurs, $S_4$ requires $q\equiv\pm1\pmod8$ and $A_5$ requires $q\equiv0,\pm1\pmod5$. This yields exactly the following table. Each entry gives $\min_{\chi\in\mathrm{Irr}(G)}\dim V_\chi^K$.
\[
\begin{array}{c|rrrrrrrrrr}
q&17&19&31&41&71&79&81&89&101&109\\\hline
K&A_4&A_4&S_4&S_4&A_5&A_5&A_5&A_5&A_5&A_5\\
\min&0&1&0&1&0&0&0&0&1&1
\end{array}
\]
For reproducibility here is the exact calculation behind the ten entries. The counts of nonidentity elements by order $(2,3,4,5)$ are $(3,8,0,0)$, $(9,8,6,0)$, and $(15,20,0,24)$, respectively. For order $n$ prime to $p$ the torus is split if $n\mid(q-1)/2$, nonsplit otherwise. All cyclic subgroups of that order are conjugate in $G$, and each contains its full set of generators in $K$ whenever it is a subgroup of $K$. For torus parameter exponent $i$, the average of its linear character over these generators is $\mu(d)/\varphi(d)$ for $d=n/\gcd(i,n)$. Thus the average on the class of elements of that order in $K$ is twice this number for a matching principal row, its negative double for a matching cuspidal row, half the corresponding value for a matching exceptional row, and zero for an opposite-series row. The trivial and Steinberg values are given in the previous section. Add the degree to the sum of the order counts times their row averages and divide by $|K|$.

The only defining-characteristic entry is $q=81$, $K=A_5$, $n=3$. All twenty such elements are conjugate in $K$, hence belong to one unipotent class of $G$. The two exceptional values there are $(1\pm9)/2$, with the sign order immaterial. The other row values are $1,-1,0$ for ordinary principal, ordinary cuspidal and Steinberg. These calculations are implemented in \texttt{check\_small\_subgroups.py} with rational arithmetic. They prove exactly the four stated overrides. As a separate check, \texttt{check\_affine.py} constructs finite fields with an explicit primitive element and counts the two unipotent types in nine candidate affine subgroups. For $q=625$ the subgroup using $r=5$ fixes dimensions $1,2$ in the exceptional pair, whereas its larger competitor using $r=25$ fixes dimensions $0,1$.

\paragraph{The remaining odd field.}
For $q=3$ the group is $A_4$. Its three linear characters are inflated from $A_4/V_4$, and its remaining irreducible is the degree-three summand of the permutation character on four letters. Their sum takes values $6,2,0$ at the identity, at involutions and at elements of order three, respectively. It is the unique minimizing character by the same positive-mean argument. A subgroup of order two fixes every linear character and one dimension in the remaining representation, since an involution there has trace $-1$. The degree lower bound $6$ rules out larger rich subgroups, so $h(3)=2$ and the optimal single-subfield degree is $6$.

\section{Further research questions}
\begin{enumerate}
\item Over the integers, does Rankin's set attain $\alpha$ as McNew conjectured in 7.2? Is $\beta>\alpha$ as asked in 7.4? The comparisons with the corresponding upper densities leave both questions open here.
\item Is the greedy ternary-valuation set optimal among sets with a natural density for \emph{rational} ratios? The finite-witness condition above reduces a negative answer to one exact weighted construction. How large is the true gap $u_2-b_2$?
\item For the cubic Artin factor, which property of the arithmetic kernel yields cubic inequalities without locating any zeros? The Gaussian perturbation calculation rules out an entire class of curvature-only proofs. A uniform argument in the Jensen degree would be needed for the full critical-line conclusion.
\item Can the optimal odd-characteristic Frobenius weights be realized in an arithmetic covolume comparison that supplies the required one-sided remainder? The character calculation gives its exact possible weights at all unramified primes.
\end{enumerate}
\end{document}
